% Generated draft: requires mathematical and source review.
\documentclass[11pt]{article}
\usepackage[margin=25mm]{geometry}
\usepackage{fontspec}
\setmainfont{Latin Modern Roman}
\usepackage{amsmath,amssymb,mathtools,unicode-math}
\setmathfont{Latin Modern Math}
\usepackage{enumitem,needspace,xurl,hyperref}
\setlength{\parindent}{0pt}
\setlength{\parskip}{5pt}
\setlength{\emergencystretch}{3em}
\newcommand{\N}{\mathbb{N}}
\newcommand{\Z}{\mathbb{Z}}
\newcommand{\Q}{\mathbb{Q}}
\newcommand{\R}{\mathbb{R}}
\newcommand{\C}{\mathbb{C}}
\newcommand{\F}{\mathbb{F}}
\newcommand{\A}{\mathbb{A}}
\newcommand{\PP}{\mathbb{P}}
\newcommand{\E}{\mathbb{E}}
\newcommand{\eps}{\varepsilon}
\newcommand{\dd}{\,\mathrm{d}}
\begin{document}
\section*{UnsolvedMath ID: 30005432}
\textbf{Source problem code:} catalogue entry
\subsection{Definitions and mathematical statement}
Use the standard left-skew-brace convention: a skew brace \(B=(B,+,\circ)\) consists of two group structures on the same set, with common identity \(0\), satisfying
\[
 a\circ(b+c)=(a\circ b)-a+(a\circ c)
 \qquad(a,b,c\in B).
\]
Neither group is assumed abelian, and \(B\) need not be finite. Write \(-a\) for the additive inverse of \(a\), and define
\[
 \lambda_a(b)=-a+(a\circ b),\qquad
 a*b=-a+(a\circ b)-b.
\]
For \(x\in B\), put
\[
 \operatorname{Fix}^{r}(x)=\{b\in B:x*b=0\},\qquad
 \operatorname{Fix}^{l}(x)=\{b\in B:b*x=0\},
\]
and define the centralizers
\[
 C_x^+=\{b\in B:b+x=x+b\},\qquad
 C_x^\circ=\{b\in B:b\circ x=x\circ b\}.
\]
An element \(x\in B\) has property \((s)\) if both subgroup indices
\[
 [(B,+):\operatorname{Fix}^{r}(x)\cap C_x^+]<\infty,
 \qquad
 [(B,\circ):\operatorname{Fix}^{l}(x)\cap C_x^\circ]<\infty
\]
are finite. Let
\[
 S(B)=\{x\in B:x\text{ has property }(s)\}.
\]
An ideal of \(B\) means a subset \(I\subseteq B\) that is a normal subgroup of both \((B,+)\) and \((B,\circ)\), and satisfies \(\lambda_a(I)\subseteq I\) for every \(a\in B\).

The question is: for every skew brace \(B\), is \(S(B)\) an ideal of \(B\)? The supplied dated assessment reports a positive result for two-sided skew braces of abelian type, meaning that \((B,+)\) is abelian and the additional identity
\[
 (a+b)\circ c=(a\circ c)-c+(b\circ c)
 \qquad(a,b,c\in B)
\]
holds. These additional assumptions are not imposed in the general question. This special-case result and the supplied open-status assessment are recorded as source claims, not independently verified facts.
\subsection{Short English statement}
In an arbitrary skew brace, consider the elements for which the two specified centralizer--fixer intersections have finite index. Must these elements form an ideal? The supplied record labels the general question open and reports a positive answer under the additional assumptions of being two-sided and having abelian additive group.
\subsection{Sources}
\begin{itemize}
\item Supplied UnsolvedMath local registry record, permanent ID \(30005432\), title \emph{The Property-s Elements of a Skew Brace}, exact source problem code \texttt{catalogue entry}, in the Oberwolfach Reports Open Problems collection. The record attributes the problem to the 2023 report and labels its statement \texttt{corrected\_verified} and its status open. Its claim that the official report was checked directly is supplied provenance, not a verification performed for this draft.
\item \emph{Mini-Workshop: Skew Braces and the Yang--Baxter Equation} (2023), Oberwolfach Reports, pp.\ 537--563, DOI 10.4171/owr/2023/9. Supplied source and literature reference labelled ``original OWR report'': \url{https://doi.org/10.4171/owr/2023/9}. The supplied original excerpt places the question in Ilaria Colazzo's contribution, \emph{Derived-indecomposable solutions and skew braces whose elements have a finite number of conjugates}, joint work with Maria Ferrara and Marco Trombetti.
\item Supplied literature reference: Colazzo--Ferrara--Trombetti, \emph{Publ. Mat.} 69 (2025): \url{https://mat.uab.es/pubmat/fitxers/download/FileType%3Apdf/FolderName%3Av69%281%29/FileName%3A6912508.pdf}.
\item Literature assessment checked at \texttt{2026-08-22}: the supplied curation states that property-\((s)\) elements form an ideal for two-sided skew braces of abelian type, while the general question remains open in the cited primary treatment. This is supplied dated curation derived from the OWR record, not independently verified literature assessment; the status and cited claims require independent verification.
\end{itemize}
\end{document}
